% ============================================================
% conclusiones.tex — Conclusiones y trabajo futuro
% ============================================================
\section{Conclusiones}

Este trabajo demostró la viabilidad de los modelos de aprendizaje automático
para la detección de fraude en transacciones financieras. Las principales
conclusiones son:

\begin{enumerate}
  \item El manejo del desbalance mediante SMOTE mejoró significativamente
        el Recall de la clase fraude respecto a los modelos sin sobremuestreo.
  \item [Completar con conclusión sobre el mejor modelo]
  \item Las variables $V_{14}$, $V_{17}$ y $V_{12}$ resultaron las más
        informativas para la detección de fraude.
  \item La solución fue desplegada como aplicación web con Flask y Plotly,
        permitiendo predicción en tiempo real.
\end{enumerate}

Como trabajo futuro se propone evaluar técnicas de detección de anomalías
no supervisadas (Isolation Forest, Autoencoder) y explorar estrategias de
monitoreo de deriva del modelo (\textit{concept drift}).

\section*{Agradecimientos}
Los autores agradecen al docente Aimer Rivera Centeno de la Universidad
Popular del Cesar por la orientación académica en este proyecto.
